% Generated draft: requires mathematical and source review.
\documentclass[11pt]{article}
\usepackage[margin=25mm]{geometry}
\usepackage{fontspec}
\setmainfont{Latin Modern Roman}
\usepackage{amsmath,amssymb,mathtools,unicode-math}
\setmathfont{Latin Modern Math}
\usepackage{enumitem,needspace,xurl,hyperref}
\setlength{\parindent}{0pt}
\setlength{\parskip}{5pt}
\setlength{\emergencystretch}{3em}
\newcommand{\N}{\mathbb{N}}
\newcommand{\Z}{\mathbb{Z}}
\newcommand{\Q}{\mathbb{Q}}
\newcommand{\R}{\mathbb{R}}
\newcommand{\C}{\mathbb{C}}
\newcommand{\F}{\mathbb{F}}
\newcommand{\A}{\mathbb{A}}
\newcommand{\PP}{\mathbb{P}}
\newcommand{\E}{\mathbb{E}}
\newcommand{\eps}{\varepsilon}
\newcommand{\dd}{\,\mathrm{d}}
\begin{document}
\section*{UnsolvedMath ID: 20003278}
\textbf{Source problem code:} catalogue entry
\subsection{Definitions and mathematical statement}
Let $G$ be a noncompact real-rank-one semisimple Lie group, $\Gamma<G$ discrete, and $U<G$ horospherical, acting on $G/\Gamma$ by left translation. F. Ledrappier and O. Sarig ask to classify nonzero $U$-invariant ergodic Radon measures, and whether each is supported on a closed $U$-orbit or comes from the harmonic-function construction. Radon means locally finite and regular; total mass need not be finite. The construction, target-group refinements and normalizations are not defined locally, so the asserted dichotomy cannot be made more specific by guessing a harmonic family.

For a standard measurable $U$-action and a nonzero $\sigma$-finite invariant ergodic measure $\mu$, let $\mathcal Z_{\mathrm{ns}}(U,\mu)$ be the invertible bimeasurable transformations commuting with $U$ (in the appropriate almost-everywhere sense) and preserving the measure class. The aid reports that every $T$ in this centralizer has a unique $c(T)>0$ with
\[
 T_*\mu=c(T)\mu,\qquad c(TS)=c(T)c(S).
\]
Thus squashability means that this character is nontrivial. Nonsingular means equivalence of measure classes, not merely preservation of null sets in one direction. The statement concerns only nonsingular centralizers, not arbitrary commuting transformations whose pushforwards might be singular.

Three reported obstructions to nontrivial scaling are: finite positive total mass; finite order of $T$ modulo null sets; or a $T$-invariant measurable set of finite positive measure. Each forces $c(T)=1$ for the relevant transformation. No conclusion that every infinite-mass measure is squashable follows.

For a normal cover with deck group $D$, the commuting deck transformations give a restriction only on the subgroup $D_\mu$ preserving the measure class. Its logarithm is a homomorphism
\[
 \chi:D_\mu\to\R,\qquad c(d)=e^{\chi(d)},
\]
factoring through $D_\mu^{\mathrm{ab}}$. If every deck transformation is nonsingular, this is a character of the full deck group and factors through $D^{\mathrm{ab}}$. Checking generators detects nontrivial deck scaling, but triviality on deck transformations does not by itself show triviality on the entire nonsingular centralizer. The deck quasi-invariance hypothesis is not supplied universally by the original question.

The aid's complete special family is $\Gamma=\{e\}$ with $U$ closed and unimodular: every nonzero ergodic $U$-invariant Radon measure on $G$ is a positive multiple of Haar measure on a single closed left-$U$ orbit $Ug$. This specifies the actual orbit, avoiding ambiguity about the source's term left coset. It is a proper free-action case, not a classification for nontrivial discrete $\Gamma$.

The source additionally reports Ledrappier--Sarig results for $u_t$ on $SL(2,\R)/\Gamma$: for a normal coabelian subgroup of a uniform lattice, the only nonsquashable ergodic measures are Haar measures up to scale; in the conilpotent case all are squashable except possibly Haar. The exact normal-cover, ergodicity and conilpotence hypotheses require their missing reference. These dated reports and the centralizer lemma do not prove the full rank-one harmonic-function dichotomy. The page footer is layout.
\subsection{Short English statement}
The aid converts squashability into a scaling character on the nonsingular centralizer and gives elementary ways that character can vanish. A deck-group check addresses only the quasi-invariant deck subgroup. The full classification is supplied only for the trivial-discrete-subgroup case, not the general horospherical problem.
\subsection{Sources}
\begin{itemize}
\item Local UnsolvedMath registry, permanent ID 20003278, source code catalogue entry; \texttt{aim-other-notes.json}, notes[84], source tag question.
\item AIM, \emph{Emerging applications of measure rigidity}, June 2004, Question 11, F. Ledrappier and O. Sarig: \url{https://aimath.org/WWN/measrigid/measrigid.pdf}.
\item Ledrappier--Sarig coabelian/conilpotent-cover squashability results and the harmonic-function construction are mentioned in the source without complete bibliographic entries or definitions.
\item Machine-generated research attempt 1: partial result, open at the research run, rigor labelled ``proved,'' low novelty confidence; registry partially solved. Not peer reviewed; independently verify before citation. Registry updated 14 August 2026.
\end{itemize}
\end{document}
